
\documentclass{article}
\usepackage{colortbl}
\usepackage{makecell}
\usepackage{multirow}
\usepackage{supertabular}

\begin{document}

\newcounter{utterance}

\centering \large Interaction Transcript for game `codenames', experiment `opponent\_difficult', episode 0 with mistral{-}large{-}3{-}azure.
\vspace{24pt}

{ \footnotesize  \setcounter{utterance}{1}
\setlength{\tabcolsep}{0pt}
\begin{supertabular}{c@{$\;$}|p{.15\linewidth}@{}p{.15\linewidth}p{.15\linewidth}p{.15\linewidth}p{.15\linewidth}p{.15\linewidth}}
    \# & $\;$A & \multicolumn{4}{c}{Game Master} & $\;\:$B\\
    \hline

    \theutterance \stepcounter{utterance}  
    & & & \multicolumn{2}{p{0.3\linewidth}}{
        \cellcolor[rgb]{0.9,0.9,0.9}{
            \makecell[{{p{\linewidth}}}]{
                \texttt{\tiny{[GM$|$GM]}}
                {'hidden': {'team': ['žyma', 'centras', 'atsargos', 'varžtas', 'traukinys', 'limuzinas', 'palydovas', 'kalnas', 'katė'], 'innocent': ['skardis', 'oras', 'siela', 'apvalus', 'vašingtonas', 'vinis', 'ligoninė'], 'opponent': ['svaras', 'piratas', 'medvilnė', 'vanduo', 'streikas', 'indija', 'viščiukas', 'amerika'], 'assassin': ['jupiteris']}, 'revealed': {'team': {'team': [], 'innocent': [], 'opponent': [], 'assassin': []}, 'opponent': {'team': [], 'innocent': [], 'opponent': [], 'assassin': []}}}
            }
        }
    }
    & & \\ \\

    \theutterance \stepcounter{utterance}  
    & & \multicolumn{4}{p{0.6\linewidth}}{
        \cellcolor[rgb]{0.9,0.9,0.9}{
            \makecell[{{p{\linewidth}}}]{
                \texttt{\tiny{[P1$\langle$GM]}}
                \texttt{Pateikite vieną vieno žodžio užuominą, susijusią su kai kuriais jūsų komandos žodžiais iš toliau pateikto sąrašo. Galite pasirinkti, su kiek ir kokiais žodžiais bus susijusi jūsų užuomina.} \\
\texttt{Tikslas {-} kad jūsų partneris, remdamasis jūsų pateiktomis užuominomis, atspėtų visus jūsų komandos žodžius, kol kita komanda neatskleidė visų savo žodžių. Jei jūsų partneris atspėja "priešininko" arba "atitraukiančio dėmesį" žodį, jo spėjimo eilė baigiasi ir priešininkų komanda gali atlikti savo ėjimą. Jei jūsų partneris atspėja "žudantį" žodį, žaidimas baigiasi ir jūsų komanda iš karto pralaimi.} \\
\texttt{Kuo daugiau žodžių, su kuriais susijusi jūsų užuomina, tuo greičiau galite laimėti žaidimą, tačiau partneriui gali būti sunkiau atspėti asociacijas, todėl išmintingai pasirinkite užuominą. Užuominos žodis turi būti semantiškai susijęs su tiksliniais žodžiais, jis negali būti vienas iš sąrašuose esančių žodžių arba sudaryti jų dalių.} \\
\\ 
\texttt{Visada pateikite vieno žodžio užuominą ir kableliais atskirtą susijusių tikslinių žodžių sąrašą toliau nurodyta forma ir atsakymus pateikite kuo trumpesnius, niekada nepateikite jokio kito teksto, nei reikalaujama šioje formoje:} \\
\\ 
\texttt{UŽUOMINA: <ŽODIS>} \\
\texttt{TIKSLAI: <TIKSLAI>} \\
\\ 
\texttt{Jūsų komandos žodžiai yra: žyma, centras, atsargos, varžtas, traukinys, limuzinas, palydovas, kalnas, katė.} \\
\texttt{Jūsų priešininko komandos žodžiai yra: svaras, piratas, medvilnė, vanduo, streikas, indija, viščiukas, amerika.} \\
\texttt{Atitraukimo žodžiai yra: skardis, oras, siela, apvalus, vašingtonas, vinis, ligoninė.} \\
\texttt{Žudymo žodžiai yra: jupiteris.} \\
            }
        }
    }
    & \\ \\

    \theutterance \stepcounter{utterance}  
    & \multicolumn{4}{p{0.6\linewidth}}{
        \cellcolor[rgb]{0.9,0.9,0.9}{
            \makecell[{{p{\linewidth}}}]{
                \texttt{\tiny{[P1$\rangle$GM]}}
                \texttt{UŽUOMINA: transportas} \\
\texttt{TIKSLAI: traukinys, limuzinas, palydovas} \\
            }
        }
    }
    & & \\ \\

    \theutterance \stepcounter{utterance}  
    & & & \multicolumn{2}{p{0.3\linewidth}}{
        \cellcolor[rgb]{0.9,0.9,0.9}{
            \makecell[{{p{\linewidth}}}]{
                \texttt{\tiny{[GM$|$GM]}}
                \texttt{transportas} \\
            }
        }
    }
    & & \\ \\

    \theutterance \stepcounter{utterance}  
    & & & \multicolumn{2}{p{0.3\linewidth}}{
        \cellcolor[rgb]{0.9,0.9,0.9}{
            \makecell[{{p{\linewidth}}}]{
                \texttt{\tiny{[GM$|$GM]}}
                ['traukinys', 'limuzinas', 'palydovas']
            }
        }
    }
    & & \\ \\

    \theutterance \stepcounter{utterance}  
    & & & \multicolumn{2}{p{0.3\linewidth}}{
        \cellcolor[rgb]{0.9,0.9,0.9}{
            \makecell[{{p{\linewidth}}}]{
                \texttt{\tiny{[GM$|$GM]}}
                {'word': 'traukinys', 'assignment': 'team'}
            }
        }
    }
    & & \\ \\

    \theutterance \stepcounter{utterance}  
    & & & \multicolumn{2}{p{0.3\linewidth}}{
        \cellcolor[rgb]{0.9,0.9,0.9}{
            \makecell[{{p{\linewidth}}}]{
                \texttt{\tiny{[GM$|$GM]}}
                {'word': 'limuzinas', 'assignment': 'team'}
            }
        }
    }
    & & \\ \\

    \theutterance \stepcounter{utterance}  
    & & & \multicolumn{2}{p{0.3\linewidth}}{
        \cellcolor[rgb]{0.9,0.9,0.9}{
            \makecell[{{p{\linewidth}}}]{
                \texttt{\tiny{[GM$|$GM]}}
                {'word': 'palydovas', 'assignment': 'team'}
            }
        }
    }
    & & \\ \\

    \theutterance \stepcounter{utterance}  
    & & \multicolumn{4}{p{0.6\linewidth}}{
        \cellcolor[rgb]{0.9,0.9,0.9}{
            \makecell[{{p{\linewidth}}}]{
                \texttt{\tiny{[GM$\rangle$P2]}}
                \texttt{Pateikite kableliais atskirtą sąrašą iki 3 žodžių iš toliau pateikto sąrašo, kurie geriausiai siejasi arba yra labiausiai susiję su žodžiu 'transportas'. Visada pradėkite savo spėjimų sąrašą žodžiais SPĖJIMAS:  ir į atsakymą neįtraukite jokio kito teksto.} \\
\\ 
\texttt{medvilnė, ligoninė, palydovas, skardis, vašingtonas, limuzinas, viščiukas, oras, kalnas, jupiteris, siela, vinis, varžtas, streikas, vanduo, traukinys, žyma, katė, indija, atsargos, piratas, amerika, svaras, apvalus, centras} \\
            }
        }
    }
    & \\ \\

    \theutterance \stepcounter{utterance}  
    & & & \multicolumn{4}{p{0.6\linewidth}}{
        \cellcolor[rgb]{0.9,0.9,0.9}{
            \makecell[{{p{\linewidth}}}]{
                \texttt{\tiny{[GM$\langle$P2]}}
                \texttt{SPĖJIMAS: traukinys, limuzinas, oras} \\
            }
        }
    }
     \\ \\

    \theutterance \stepcounter{utterance}  
    & & & \multicolumn{2}{p{0.3\linewidth}}{
        \cellcolor[rgb]{0.9,0.9,0.9}{
            \makecell[{{p{\linewidth}}}]{
                \texttt{\tiny{[GM$|$GM]}}
                ['traukinys', 'limuzinas', 'oras']
            }
        }
    }
    & & \\ \\

    \theutterance \stepcounter{utterance}  
    & & & \multicolumn{2}{p{0.3\linewidth}}{
        \cellcolor[rgb]{0.9,0.9,0.9}{
            \makecell[{{p{\linewidth}}}]{
                \texttt{\tiny{[GM$|$GM]}}
                {'word': 'traukinys', 'assignment': 'team'}
            }
        }
    }
    & & \\ \\

    \theutterance \stepcounter{utterance}  
    & & & \multicolumn{2}{p{0.3\linewidth}}{
        \cellcolor[rgb]{0.9,0.9,0.9}{
            \makecell[{{p{\linewidth}}}]{
                \texttt{\tiny{[GM$|$GM]}}
                {'word': 'traukinys', 'assignment': 'team'}
            }
        }
    }
    & & \\ \\

    \theutterance \stepcounter{utterance}  
    & & & \multicolumn{2}{p{0.3\linewidth}}{
        \cellcolor[rgb]{0.9,0.9,0.9}{
            \makecell[{{p{\linewidth}}}]{
                \texttt{\tiny{[GM$|$GM]}}
                {'word': 'limuzinas', 'assignment': 'team'}
            }
        }
    }
    & & \\ \\

    \theutterance \stepcounter{utterance}  
    & & & \multicolumn{2}{p{0.3\linewidth}}{
        \cellcolor[rgb]{0.9,0.9,0.9}{
            \makecell[{{p{\linewidth}}}]{
                \texttt{\tiny{[GM$|$GM]}}
                {'word': 'limuzinas', 'assignment': 'team'}
            }
        }
    }
    & & \\ \\

    \theutterance \stepcounter{utterance}  
    & & & \multicolumn{2}{p{0.3\linewidth}}{
        \cellcolor[rgb]{0.9,0.9,0.9}{
            \makecell[{{p{\linewidth}}}]{
                \texttt{\tiny{[GM$|$GM]}}
                {'word': 'oras', 'assignment': 'innocent'}
            }
        }
    }
    & & \\ \\

    \theutterance \stepcounter{utterance}  
    & & & \multicolumn{2}{p{0.3\linewidth}}{
        \cellcolor[rgb]{0.9,0.9,0.9}{
            \makecell[{{p{\linewidth}}}]{
                \texttt{\tiny{[GM$|$GM]}}
                \texttt{oras} \\
            }
        }
    }
    & & \\ \\

    \theutterance \stepcounter{utterance}  
    & & & \multicolumn{2}{p{0.3\linewidth}}{
        \cellcolor[rgb]{0.9,0.9,0.9}{
            \makecell[{{p{\linewidth}}}]{
                \texttt{\tiny{[GM$|$GM]}}
                {'word': 'streikas', 'assignment': 'opponent'}
            }
        }
    }
    & & \\ \\

    \theutterance \stepcounter{utterance}  
    & & & \multicolumn{2}{p{0.3\linewidth}}{
        \cellcolor[rgb]{0.9,0.9,0.9}{
            \makecell[{{p{\linewidth}}}]{
                \texttt{\tiny{[GM$|$GM]}}
                {'word': 'viščiukas', 'assignment': 'opponent'}
            }
        }
    }
    & & \\ \\

    \theutterance \stepcounter{utterance}  
    & & & \multicolumn{2}{p{0.3\linewidth}}{
        \cellcolor[rgb]{0.9,0.9,0.9}{
            \makecell[{{p{\linewidth}}}]{
                \texttt{\tiny{[GM$|$GM]}}
                {'hidden': {'team': ['žyma', 'centras', 'atsargos', 'varžtas', 'palydovas', 'kalnas', 'katė'], 'innocent': ['skardis', 'siela', 'apvalus', 'vašingtonas', 'vinis', 'ligoninė'], 'opponent': ['svaras', 'piratas', 'medvilnė', 'vanduo', 'indija', 'amerika'], 'assassin': ['jupiteris']}, 'revealed': {'team': {'team': ['traukinys', 'limuzinas'], 'innocent': ['oras'], 'opponent': [], 'assassin': []}, 'opponent': {'team': [], 'innocent': [], 'opponent': ['streikas', 'viščiukas'], 'assassin': []}}}
            }
        }
    }
    & & \\ \\

    \theutterance \stepcounter{utterance}  
    & & \multicolumn{4}{p{0.6\linewidth}}{
        \cellcolor[rgb]{0.9,0.9,0.9}{
            \makecell[{{p{\linewidth}}}]{
                \texttt{\tiny{[P1$\langle$GM]}}
                \texttt{Žodžiai traukinys, limuzinas buvo atspėti teisingai. Žodis oras buvo atspėtas, tačiau tai nekaltas žodis. Jūsų komandos draugo ėjimas baigėsi.} \\
\texttt{Dabar pateikite kitą užuominą, susijusią su kai kuriais likusiais jūsų komandos žodžiais, ir susijusių tikslinių žodžių sąrašą. Nepamirškite užuominą pradėti žodžiais "UŽUOMINA: ", įrašykite naują eilutę, o tikslinių žodžių sąrašą, atskirtą kableliais, pradėkite žodžiais "TIKSLAI: ". Pastaba: palyginti su ankstesniais prašymais, kai kurie žodžiai iš sąrašų buvo pašalinti.} \\
\\ 
\texttt{Likę jūsų komandos žodžiai yra šie: žyma, centras, atsargos, varžtas, palydovas, kalnas, katė.} \\
\texttt{Likę priešininko žodžiai yra šie: svaras, piratas, medvilnė, vanduo, streikas, indija, viščiukas, amerika.} \\
\texttt{Likę atitraukiantys žodžiai yra: skardis, siela, apvalus, vašingtonas, vinis, ligoninė.} \\
\texttt{Likę nužudymo žodžiai yra: jupiteris.} \\
            }
        }
    }
    & \\ \\

    \theutterance \stepcounter{utterance}  
    & \multicolumn{4}{p{0.6\linewidth}}{
        \cellcolor[rgb]{0.9,0.9,0.9}{
            \makecell[{{p{\linewidth}}}]{
                \texttt{\tiny{[P1$\rangle$GM]}}
                \texttt{UŽUOMINA: tvirtinimas} \\
\texttt{TIKSLAI: varžtas, atsargos} \\
            }
        }
    }
    & & \\ \\

    \theutterance \stepcounter{utterance}  
    & & & \multicolumn{2}{p{0.3\linewidth}}{
        \cellcolor[rgb]{0.9,0.9,0.9}{
            \makecell[{{p{\linewidth}}}]{
                \texttt{\tiny{[GM$|$GM]}}
                \texttt{tvirtinimas} \\
            }
        }
    }
    & & \\ \\

    \theutterance \stepcounter{utterance}  
    & & & \multicolumn{2}{p{0.3\linewidth}}{
        \cellcolor[rgb]{0.9,0.9,0.9}{
            \makecell[{{p{\linewidth}}}]{
                \texttt{\tiny{[GM$|$GM]}}
                ['varžtas', 'atsargos']
            }
        }
    }
    & & \\ \\

    \theutterance \stepcounter{utterance}  
    & & & \multicolumn{2}{p{0.3\linewidth}}{
        \cellcolor[rgb]{0.9,0.9,0.9}{
            \makecell[{{p{\linewidth}}}]{
                \texttt{\tiny{[GM$|$GM]}}
                {'word': 'varžtas', 'assignment': 'team'}
            }
        }
    }
    & & \\ \\

    \theutterance \stepcounter{utterance}  
    & & & \multicolumn{2}{p{0.3\linewidth}}{
        \cellcolor[rgb]{0.9,0.9,0.9}{
            \makecell[{{p{\linewidth}}}]{
                \texttt{\tiny{[GM$|$GM]}}
                {'word': 'atsargos', 'assignment': 'team'}
            }
        }
    }
    & & \\ \\

    \theutterance \stepcounter{utterance}  
    & & \multicolumn{4}{p{0.6\linewidth}}{
        \cellcolor[rgb]{0.9,0.9,0.9}{
            \makecell[{{p{\linewidth}}}]{
                \texttt{\tiny{[GM$\rangle$P2]}}
                \texttt{Žodžiai traukinys, limuzinas buvo atspėti teisingai. Žodis oras buvo atspėtas, tačiau tai nekaltas žodis. Jūsų ėjimas baigėsi.} \\
\texttt{Dabar pateikite dar vieną kableliais atskirtą sąrašą iš bent 1 ir ne daugiau kaip 2 žodžių iš toliau pateikto žodžių sąrašo, kurie geriausiai siejasi arba yra labiausiai susiję su žodžiu 'tvirtinimas'. Nepamirškite savo atsakymą pradėti SPĖJIMAS: . Pastaba: palyginti su ankstesniais prašymais, kai kurie žodžiai iš sąrašo buvo pašalinti.} \\
\\ 
\texttt{medvilnė, ligoninė, palydovas, skardis, vašingtonas, kalnas, jupiteris, siela, vinis, varžtas, vanduo, žyma, katė, indija, atsargos, piratas, amerika, svaras, apvalus, centras} \\
            }
        }
    }
    & \\ \\

    \theutterance \stepcounter{utterance}  
    & & & \multicolumn{4}{p{0.6\linewidth}}{
        \cellcolor[rgb]{0.9,0.9,0.9}{
            \makecell[{{p{\linewidth}}}]{
                \texttt{\tiny{[GM$\langle$P2]}}
                \texttt{SPĖJIMAS: vinis, varžtas} \\
            }
        }
    }
     \\ \\

    \theutterance \stepcounter{utterance}  
    & & & \multicolumn{2}{p{0.3\linewidth}}{
        \cellcolor[rgb]{0.9,0.9,0.9}{
            \makecell[{{p{\linewidth}}}]{
                \texttt{\tiny{[GM$|$GM]}}
                ['vinis', 'varžtas']
            }
        }
    }
    & & \\ \\

    \theutterance \stepcounter{utterance}  
    & & & \multicolumn{2}{p{0.3\linewidth}}{
        \cellcolor[rgb]{0.9,0.9,0.9}{
            \makecell[{{p{\linewidth}}}]{
                \texttt{\tiny{[GM$|$GM]}}
                {'word': 'vinis', 'assignment': 'innocent'}
            }
        }
    }
    & & \\ \\

    \theutterance \stepcounter{utterance}  
    & & & \multicolumn{2}{p{0.3\linewidth}}{
        \cellcolor[rgb]{0.9,0.9,0.9}{
            \makecell[{{p{\linewidth}}}]{
                \texttt{\tiny{[GM$|$GM]}}
                \texttt{vinis} \\
            }
        }
    }
    & & \\ \\

    \theutterance \stepcounter{utterance}  
    & & & \multicolumn{2}{p{0.3\linewidth}}{
        \cellcolor[rgb]{0.9,0.9,0.9}{
            \makecell[{{p{\linewidth}}}]{
                \texttt{\tiny{[GM$|$GM]}}
                {'word': 'svaras', 'assignment': 'opponent'}
            }
        }
    }
    & & \\ \\

    \theutterance \stepcounter{utterance}  
    & & & \multicolumn{2}{p{0.3\linewidth}}{
        \cellcolor[rgb]{0.9,0.9,0.9}{
            \makecell[{{p{\linewidth}}}]{
                \texttt{\tiny{[GM$|$GM]}}
                {'word': 'piratas', 'assignment': 'opponent'}
            }
        }
    }
    & & \\ \\

    \theutterance \stepcounter{utterance}  
    & & & \multicolumn{2}{p{0.3\linewidth}}{
        \cellcolor[rgb]{0.9,0.9,0.9}{
            \makecell[{{p{\linewidth}}}]{
                \texttt{\tiny{[GM$|$GM]}}
                {'hidden': {'team': ['žyma', 'centras', 'atsargos', 'varžtas', 'palydovas', 'kalnas', 'katė'], 'innocent': ['skardis', 'siela', 'apvalus', 'vašingtonas', 'ligoninė'], 'opponent': ['medvilnė', 'vanduo', 'indija', 'amerika'], 'assassin': ['jupiteris']}, 'revealed': {'team': {'team': ['traukinys', 'limuzinas'], 'innocent': ['oras', 'vinis'], 'opponent': [], 'assassin': []}, 'opponent': {'team': [], 'innocent': [], 'opponent': ['streikas', 'viščiukas', 'svaras', 'piratas'], 'assassin': []}}}
            }
        }
    }
    & & \\ \\

    \theutterance \stepcounter{utterance}  
    & & \multicolumn{4}{p{0.6\linewidth}}{
        \cellcolor[rgb]{0.9,0.9,0.9}{
            \makecell[{{p{\linewidth}}}]{
                \texttt{\tiny{[P1$\langle$GM]}}
                \texttt{Žodis vinis buvo atspėtas, tačiau tai nekaltas žodis. Jūsų komandos draugo ėjimas baigėsi.} \\
\texttt{Dabar pateikite kitą užuominą, susijusią su kai kuriais likusiais jūsų komandos žodžiais, ir susijusių tikslinių žodžių sąrašą. Nepamirškite užuominą pradėti žodžiais "UŽUOMINA: ", įrašykite naują eilutę, o tikslinių žodžių sąrašą, atskirtą kableliais, pradėkite žodžiais "TIKSLAI: ". Pastaba: palyginti su ankstesniais prašymais, kai kurie žodžiai iš sąrašų buvo pašalinti.} \\
\\ 
\texttt{Likę jūsų komandos žodžiai yra šie: žyma, centras, atsargos, varžtas, palydovas, kalnas, katė.} \\
\texttt{Likę priešininko žodžiai yra šie: svaras, piratas, medvilnė, vanduo, indija, amerika.} \\
\texttt{Likę atitraukiantys žodžiai yra: skardis, siela, apvalus, vašingtonas, ligoninė.} \\
\texttt{Likę nužudymo žodžiai yra: jupiteris.} \\
            }
        }
    }
    & \\ \\

    \theutterance \stepcounter{utterance}  
    & \multicolumn{4}{p{0.6\linewidth}}{
        \cellcolor[rgb]{0.9,0.9,0.9}{
            \makecell[{{p{\linewidth}}}]{
                \texttt{\tiny{[P1$\rangle$GM]}}
                \texttt{UŽUOMINA: žvaigždė} \\
\texttt{TIKSLAI: žyma, palydovas} \\
            }
        }
    }
    & & \\ \\

    \theutterance \stepcounter{utterance}  
    & & & \multicolumn{2}{p{0.3\linewidth}}{
        \cellcolor[rgb]{0.9,0.9,0.9}{
            \makecell[{{p{\linewidth}}}]{
                \texttt{\tiny{[GM$|$GM]}}
                \texttt{žvaigždė} \\
            }
        }
    }
    & & \\ \\

    \theutterance \stepcounter{utterance}  
    & & & \multicolumn{2}{p{0.3\linewidth}}{
        \cellcolor[rgb]{0.9,0.9,0.9}{
            \makecell[{{p{\linewidth}}}]{
                \texttt{\tiny{[GM$|$GM]}}
                ['žyma', 'palydovas']
            }
        }
    }
    & & \\ \\

    \theutterance \stepcounter{utterance}  
    & & & \multicolumn{2}{p{0.3\linewidth}}{
        \cellcolor[rgb]{0.9,0.9,0.9}{
            \makecell[{{p{\linewidth}}}]{
                \texttt{\tiny{[GM$|$GM]}}
                {'word': 'žyma', 'assignment': 'team'}
            }
        }
    }
    & & \\ \\

    \theutterance \stepcounter{utterance}  
    & & & \multicolumn{2}{p{0.3\linewidth}}{
        \cellcolor[rgb]{0.9,0.9,0.9}{
            \makecell[{{p{\linewidth}}}]{
                \texttt{\tiny{[GM$|$GM]}}
                {'word': 'palydovas', 'assignment': 'team'}
            }
        }
    }
    & & \\ \\

    \theutterance \stepcounter{utterance}  
    & & \multicolumn{4}{p{0.6\linewidth}}{
        \cellcolor[rgb]{0.9,0.9,0.9}{
            \makecell[{{p{\linewidth}}}]{
                \texttt{\tiny{[GM$\rangle$P2]}}
                \texttt{Žodis vinis buvo atspėtas, tačiau tai nekaltas žodis. Jūsų ėjimas baigėsi.} \\
\texttt{Dabar pateikite dar vieną kableliais atskirtą sąrašą iš bent 1 ir ne daugiau kaip 2 žodžių iš toliau pateikto žodžių sąrašo, kurie geriausiai siejasi arba yra labiausiai susiję su žodžiu 'žvaigždė'. Nepamirškite savo atsakymą pradėti SPĖJIMAS: . Pastaba: palyginti su ankstesniais prašymais, kai kurie žodžiai iš sąrašo buvo pašalinti.} \\
\\ 
\texttt{medvilnė, ligoninė, palydovas, skardis, vašingtonas, kalnas, jupiteris, siela, varžtas, vanduo, žyma, katė, indija, atsargos, amerika, apvalus, centras} \\
            }
        }
    }
    & \\ \\

    \theutterance \stepcounter{utterance}  
    & & & \multicolumn{4}{p{0.6\linewidth}}{
        \cellcolor[rgb]{0.9,0.9,0.9}{
            \makecell[{{p{\linewidth}}}]{
                \texttt{\tiny{[GM$\langle$P2]}}
                \texttt{SPĖJIMAS: palydovas, jupiteris} \\
            }
        }
    }
     \\ \\

    \theutterance \stepcounter{utterance}  
    & & & \multicolumn{2}{p{0.3\linewidth}}{
        \cellcolor[rgb]{0.9,0.9,0.9}{
            \makecell[{{p{\linewidth}}}]{
                \texttt{\tiny{[GM$|$GM]}}
                ['palydovas', 'jupiteris']
            }
        }
    }
    & & \\ \\

    \theutterance \stepcounter{utterance}  
    & & & \multicolumn{2}{p{0.3\linewidth}}{
        \cellcolor[rgb]{0.9,0.9,0.9}{
            \makecell[{{p{\linewidth}}}]{
                \texttt{\tiny{[GM$|$GM]}}
                {'word': 'palydovas', 'assignment': 'team'}
            }
        }
    }
    & & \\ \\

    \theutterance \stepcounter{utterance}  
    & & & \multicolumn{2}{p{0.3\linewidth}}{
        \cellcolor[rgb]{0.9,0.9,0.9}{
            \makecell[{{p{\linewidth}}}]{
                \texttt{\tiny{[GM$|$GM]}}
                {'word': 'palydovas', 'assignment': 'team'}
            }
        }
    }
    & & \\ \\

    \theutterance \stepcounter{utterance}  
    & & & \multicolumn{2}{p{0.3\linewidth}}{
        \cellcolor[rgb]{0.9,0.9,0.9}{
            \makecell[{{p{\linewidth}}}]{
                \texttt{\tiny{[GM$|$GM]}}
                {'word': 'jupiteris', 'assignment': 'assassin'}
            }
        }
    }
    & & \\ \\

    \theutterance \stepcounter{utterance}  
    & & & \multicolumn{2}{p{0.3\linewidth}}{
        \cellcolor[rgb]{0.9,0.9,0.9}{
            \makecell[{{p{\linewidth}}}]{
                \texttt{\tiny{[GM$|$GM]}}
                \texttt{jupiteris} \\
            }
        }
    }
    & & \\ \\

    \theutterance \stepcounter{utterance}  
    & & & \multicolumn{2}{p{0.3\linewidth}}{
        \cellcolor[rgb]{0.9,0.9,0.9}{
            \makecell[{{p{\linewidth}}}]{
                \texttt{\tiny{[GM$|$GM]}}
                \texttt{opponent has won through assassin} \\
            }
        }
    }
    & & \\ \\

\end{supertabular}
}

\end{document}
